\subsection{Desktop-Integration mit Tauri und Rust}
\label{subsec:tauri-rust-desktop}

Tauri bildet eine eigene native Grenze, ohne das Frontend zu duplizieren.
\path{src-tauri/tauri.conf.json} startet für die Entwicklung ausschließlich den
Vite-Server und verwendet für Builds \path{frontend/dist}. Der Rust-Einstieg
erstellt lediglich die Tauri-Anwendung; native Fachkommandos oder ein
eingebetteter FastAPI-Sidecar sind im Template nicht implementiert. Dies nutzt
Tauris Fähigkeit, Web-Frontends in einer nativen Webview zu hosten, während
Rust die native Kompositionsgrenze bereitstellt
\parencite{tauri2026start,rust2026book,kleiveist2026architecture}.

Die Standard-Capability gewährt nur \texttt{core:default}; eine Content Security
Policy begrenzt lokale Ressourcen und die dokumentierten Entwicklungsendpunkte.
Weitere native Berechtigungen müssen produktspezifisch und möglichst eng
ergänzt werden \parencite{tauri2026capabilities}. Im Profil Desktop-local bleibt
Funktionalität lokal im Frontend oder in später ergänzten Tauri-Kommandos.
Desktop-cloud und Full-platform verwenden dagegen die öffentliche FastAPI-URL.
Packaging, Signierung und die Entscheidung zwischen Remote-API, Sidecar oder
lokalen Kommandos sind Produktentscheidungen und nicht Teil dieser Baseline.
